% Self-contained contribution for the Polylogarithms continuation article.
% Required packages: amsmath, amssymb, amsthm, booktabs, hyperref.
% Required theorem environments: theorem, lemma, corollary, remark.

\section{Proof certificates for the cubic and quartic trilogarithm ladders}
\label{gaussian:sec:proved-small-ladders}

The two displayed formulas record two trilogarithm combinations
at the positive roots of $r^3+r=1$ and $r^4+r=1$, with substantial experimental evidence. Their exact functional-equation
certificates now explain the coefficients.
We give such derivations. The quartic identity belongs to a uniform
family; the cubic identity has a short certificate using two
specializations of Kummer's nine-term relation and four Rogers pentagons.
These are exact consequences of classical functional equations, not
claims that the underlying equations are new.

\subsection{A uniform ladder from a complementary pair of powers}

For $0<x<1$ the real Landen relation for the trilogarithm is
\begin{align}
\operatorname{Li}_3(x)+\operatorname{Li}_3(1-x)
 +\operatorname{Li}_3(1-x^{-1})
 &=\zeta(3)+\frac{\pi^2}{6}\log x
 -\frac12\log^2x\log(1-x)+\frac16\log^3x.
\label{gaussian:eq:real-trilog-landen}
\end{align}
For a modern account and its arithmetic analogues see
Nakamura and Shiraishi, \emph{Landen's trilogarithm functional equation
and $\ell$-adic Galois multiple polylogarithms} \cite{gaussian:NakamuraShiraishi}.
For completeness, a real-variable proof is especially short.
Differentiation of the left side, followed by
\[
\operatorname{Li}_2(x)+\operatorname{Li}_2(1-x)
 =\frac{\pi^2}{6}-\log x\log(1-x),\qquad
\operatorname{Li}_2(1-x^{-1})
 =-\operatorname{Li}_2(1-x)-\frac12\log^2x,
\]
gives
\[
\frac{\pi^2/6-\log x\log(1-x)}{x}
  +\frac{\log^2x}{2x(1-x)}.
\]
This is the derivative of the right side. The difference tends to zero
as $x\to1^-$, proving \eqref{gaussian:eq:real-trilog-landen}. The two dilogarithm
relations used here follow themselves by one differentiation and an
endpoint value. All logarithms in the calculation are real, so there
is no branch ambiguity.

\begin{theorem}[Complementary-power trilogarithm ladder]
\label{gaussian:thm:complementary-powers}
Let $a,b$ be positive integers with $a\le b$, and let $0<\rho<1$ be
the unique positive solution of $\rho^a+\rho^b=1$. Set
$\lambda=\log\rho$ and $d=b-a$. Then
\begin{align}
&\operatorname{Li}_3(\rho^a)+\operatorname{Li}_3(\rho^b)
 -\operatorname{Li}_3(\rho^d)
 +\frac14\operatorname{Li}_3(\rho^{2d})\notag\\
&\hspace{20mm}=\zeta(3)+\frac{a\pi^2}{6}\lambda
 +\frac{a^2(a-3b)}{6}\lambda^3.
\label{gaussian:eq:complementary-powers}
\end{align}
When $a=b$, the powers $\rho^0$ are interpreted as $1$.
\end{theorem}

\begin{proof}
The polynomial $t^a+t^b$ is strictly increasing from $0$ to $2$ on
$[0,1]$, which proves existence and uniqueness. Substitute $x=\rho^a$
in \eqref{gaussian:eq:real-trilog-landen}; its other arguments are
$1-x=\rho^b$ and $1-x^{-1}=-\rho^{b-a}$. The absolutely convergent
series identity
\[
\operatorname{Li}_3(z)+\operatorname{Li}_3(-z)
 =\frac14\operatorname{Li}_3(z^2),\qquad 0\le z\le1,
\]
eliminates the negative argument. The logarithmic terms become exactly
the right side of \eqref{gaussian:eq:complementary-powers}.
\end{proof}

In particular, if $r_m^m+r_m=1$ and $\lambda_m=\log r_m$, then
\begin{align}
\operatorname{Li}_3(r_m)+\operatorname{Li}_3(r_m^m)
 -\operatorname{Li}_3(r_m^{m-1})
 +\frac14\operatorname{Li}_3(r_m^{2m-2})
 =\zeta(3)+\frac{\pi^2}{6}\lambda_m
 -\frac{3m-1}{6}\lambda_m^3.
\label{gaussian:eq:all-m-trilog}
\end{align}
This supplies an infinite family of finite exact reductions, although
different pairs $(a,b)$ can describe the same base after a power
substitution. It does not assert that these exhaust all ladders for
any one number field.

\begin{corollary}[The quartic identity]
If $r^4+r=1$ and $\lambda=\log r$, then
\begin{align}
-12\operatorname{Li}_3(r)+12\operatorname{Li}_3(r^3)
 -12\operatorname{Li}_3(r^4)-3\operatorname{Li}_3(r^6)
 =22\lambda^3-2\pi^2\lambda-12\zeta(3).
\label{gaussian:eq:quartic-proved}
\end{align}
\end{corollary}
\begin{proof}
Take $(a,b)=(1,4)$ in Theorem~\ref{gaussian:thm:complementary-powers}
and multiply by $-12$.
\end{proof}

For $m=2$, \eqref{gaussian:eq:all-m-trilog} reduces to the familiar golden value
\[
\operatorname{Li}_3(r_2^2)
 =\frac45\zeta(3)+\frac{2\pi^2}{15}\log r_2
 -\frac23\log^3r_2.
\]
For $m=3$ it gives a useful additional relation on the arguments
$r,r^2,r^3,r^4$; the manuscript's cubic identity instead uses
$r,r^2,r^3,r^6$. The change of the last argument requires more than
the one-variable Landen relation.

\subsection{An elementary dilogarithm certificate in the cubic field}

Throughout this subsection let $r^3+r=1$, $0<r<1$, and put
$\lambda=\log r$, $A=\log(1+r)$. Direct algebra in
$\mathbb Q[r]/(r^3+r-1)$ gives
\begin{align}
1-r&=r^3,&1-r^2&=r^3(1+r),&1-r^3&=r,\\
1-r^4&=r^2(1+r),&1+r^3&=r^2(1+r)^2,
 &1-r^6&=r^3(1+r)^2.
\label{gaussian:eq:cubic-factorizations}
\end{align}
Define the Rogers dilogarithm on $(0,1)$ by
\[
R(x)=\operatorname{Li}_2(x)+\frac12\log x\log(1-x).
\]
It satisfies $R(x)+R(1-x)=\zeta(2)$ and the pentagon
\begin{align}
R(x)+R(y)
 =R(xy)+R\!\left(\frac{x(1-y)}{1-xy}\right)
  +R\!\left(\frac{y(1-x)}{1-xy}\right),\qquad 0<x,y<1.
\label{gaussian:eq:rogers-pentagon}
\end{align}
The latter is a real-domain form of the classical five-term relation.
It may also be verified by differentiating in $x$ with fixed $y$ and
using the limit at $x=0$.

\begin{lemma}[Cubic dilogarithm identity]
\label{gaussian:lem:cubic-dilog}
\begin{equation}
\operatorname{Li}_2(r^6)
 =6\operatorname{Li}_2(r^2)-2\operatorname{Li}_2(r)
  -4\operatorname{Li}_2(r^3).
\label{gaussian:eq:cubic-dilog}
\end{equation}
\end{lemma}

\begin{proof}
Set $s=r/(1+r)$ and $t=r^4/(1+r)$, so every argument below lies
in $(0,1)$. Write $R_k=R(r^k)$ and $R_s=R(s)$, $R_t=R(t)$.
Applying \eqref{gaussian:eq:rogers-pentagon} at $(x,y)=(r,r)$ gives
\[
2R_1=R_2+2R_s.
\]
At $(x,y)=(r,r^3)$ its two remaining arguments are $1-s$ and $t$.
Since $R_1+R_3=\zeta(2)$, reflection gives
\[
R_t=R_s-R_4.
\]
At $(x,y)=(r^2,r^4)$ its two remaining arguments are $s$ and $t$,
and therefore
\[
R_6=R_2+2R_4-2R_s.
\]
Finally, at $(x,y)=(r,r^2)$ these arguments are $1-r^2$ and $r^4$.
Using reflection twice gives
\[
R_4=2R_2-2R_3.
\]
Eliminating $R_s,R_t,R_4$ now yields
$R_6=6R_2-2R_1-4R_3$. The difference between $R(r^k)$ and
$\operatorname{Li}_2(r^k)$ is
$(k\lambda/2)\log(1-r^k)$. In the indicated combination these
terms cancel: by \eqref{gaussian:eq:cubic-factorizations}, the required check is
\[
3\lambda(3\lambda+2A)
 -6\lambda(3\lambda+A)+3\lambda^2+6\lambda^2=0.
\]
This proves the claimed ordinary dilogarithm identity.
\end{proof}

\subsection{Two Kummer substitutions prove the cubic trilogarithm}

The single-valued trilogarithm used in the manuscript is
\[
P_3(z)=\operatorname{Re}\operatorname{Li}_3(z)
 -\log|z|\operatorname{Re}\operatorname{Li}_2(z)
 -\frac13\log^2|z|\log|1-z|.
\]
The value at $z=1$ is defined by continuity. It obeys
\[
P_3(z^{-1})=P_3(z),\qquad
P_3(z)+P_3(-z)=\frac14P_3(z^2),
\]
and the single-valued Landen relation
$P_3(x)+P_3(1-x)+P_3(1-x^{-1})=\zeta(3)$ for real $0<x<1$.
The latter follows from \eqref{gaussian:eq:real-trilog-landen} and the two
dilogarithm identities in its proof.

We use Kummer's classical equation in the form stated and justified
by Zagier, \emph{Special values and functional equations of
polylogarithms}, \S7, Example~3, equation~(7) \cite{gaussian:Zagier}:
\begin{align}
&2P_3(x)+2P_3(y)
 +2P_3\!\left(\frac{x(1-y)}{x-1}\right)
 +2P_3\!\left(\frac{y(1-x)}{y-1}\right)\notag\\
&\quad+2P_3\!\left(\frac{1-x}{1-y}\right)
 +2P_3\!\left(\frac{x(1-y)}{y(1-x)}\right)
 -P_3(xy)-P_3(x/y)\notag\\
&\quad-P_3\!\left(\frac{x(1-y)^2}{y(1-x)^2}\right)=2\zeta(3).
\label{gaussian:eq:kummer-P3}
\end{align}
We only use specializations with nonzero denominators and finite
arguments. Since $P_3$ is single-valued, some real arguments exceeding
$1$ cause no branch correction in this equation; after inversion and
duplication, all remaining arguments are $r^k\in(0,1)$.

\begin{theorem}[The cubic identity]
\label{gaussian:thm:cubic-proved}
Let $r^3+r=1$ with $0<r<1$, and let $\lambda=\log r$. Then
\begin{align}
&-60\operatorname{Li}_3(r)+54\operatorname{Li}_3(r^2)
 -32\operatorname{Li}_3(r^3)-3\operatorname{Li}_3(r^6)\notag\\
&\hspace{20mm}=24\lambda^3-4\pi^2\lambda-36\zeta(3).
\label{gaussian:eq:cubic-proved}
\end{align}
\end{theorem}

\begin{proof}
Write $P_k=P_3(r^k)$. Landen at $r$ and duplication give
\[
\mathsf A:=P_1-P_2+P_3+\tfrac14P_4=\zeta(3).
\]
The nine arguments of \eqref{gaussian:eq:kummer-P3} at
$(x,y)=(r,r^3)$ and $(-r,r^2)$, in their displayed order, are
\[
\begin{array}{c|rrrrrrrrr}
 &1&2&3&4&5&6&7&8&9\\ \hline
(r,r^3)&r&r^3&-r^{-1}&-r^5&r^2&r^{-4}&r^4&r^{-2}&r^{-6}\\
(-r,r^2)&-r&r^2&r^4&-r^{-1}&r^{-3}&-r^2&-r^3&-r^{-1}&-r^5
\end{array}
\]
respectively. Using inversion and
$P_3(-r^k)=\tfrac14P_{2|k|}-P_{|k|}$, these two equations reduce to
\begin{align*}
\mathsf B&:=\tfrac32P_2+2P_3+P_4-2P_5-P_6+\tfrac12P_{10}
 =2\zeta(3),\\
\mathsf C&:=-3P_1+\tfrac34P_2+3P_3+\tfrac52P_4+P_5
 -\tfrac14P_6-\tfrac14P_{10}=2\zeta(3).
\end{align*}
The exact rational row operation $4\mathsf C+2\mathsf B-48\mathsf A$
therefore proves
\begin{equation}
-60P_1+54P_2-32P_3-3P_6=-36\zeta(3).
\label{gaussian:eq:cubic-P3-certificate}
\end{equation}

To recover ordinary trilogarithms, put
$c_1=-60,c_2=54,c_3=-32,c_6=-3$ and
\[
D=\sum_{k\in\{1,2,3,6\}}c_k\operatorname{Li}_3(r^k),\quad
S=\sum_{k\in\{1,2,3,6\}}c_k k\operatorname{Li}_2(r^k),\quad
T=\sum_{k\in\{1,2,3,6\}}c_k k^2\log(1-r^k).
\]
Expansion of the definition gives
$\sum c_kP_k=D-\lambda S-\lambda^2 T/3$.
Lemma~\ref{gaussian:lem:cubic-dilog}, together with
$\operatorname{Li}_2(r)+\operatorname{Li}_2(r^3)
=\pi^2/6-3\lambda^2$, gives
\[
S=72\lambda^2-4\pi^2.
\]
The factorizations \eqref{gaussian:eq:cubic-factorizations} give
\[
T=-60(3\lambda)+216(3\lambda+A)
 -288\lambda-108(3\lambda+2A)=-144\lambda.
\]
Consequently \eqref{gaussian:eq:cubic-P3-certificate} becomes
$D-24\lambda^3+4\pi^2\lambda=-36\zeta(3)$, which is
\eqref{gaussian:eq:cubic-proved}.
\end{proof}

\subsection{What the certificate establishes and what remains open}

The proofs settle the two displayed cubic and quartic trilogarithm
identities without period-independence assumptions, numerical recognition,
or unproved higher Bloch-group assertions. The universal family
\eqref{gaussian:eq:complementary-powers} also gives a principled generator of
additional exact identities. It does not justify pseudointegration to
weight four or above: differentiating with respect to an algebraic base
does not preserve its isolated defining relation.

Natural next problems are to derive comparable local certificates for the
weight-five through weight-nine golden formulas, find the smallest useful
functional-equation closure for a chosen base, and extend the exact
argument/substitution checker to arbitrary algebraic number fields.
Any claim of minimality of the surviving period basket remains distinct
from the existence of the displayed reductions.
